\section{Background}
%\vspace{-3mm}
\textbf{Markov Decision Processes:}
Sequential decision making tasks can be modeled by
Markov Decision Processes (MDPs)~\citep{sutton2018reinforcement}~%~\cite{puterman1990markov,sutton2018reinforcement}
represented by a tuple $\gM = \big( \gS, \gA, R, P, \gamma, \mu \big)$, where $\gS$ is the state space, $\gA$ is the action space, $R: \gS \times \gA \rightarrow [R_{\min}, R_{\max} ]$ is the reward function
% \footnote{We assume the reward to be deterministic to simplify notations but it can be stochastic in general.}
, $P: \gS \times \gA \times \gS \rightarrow \sR_{\geq 0}$ is the transition function, $\gamma \in (0, 1]$ is the discount factor, and $\mu: \gS \rightarrow \sR_{\geq 0}$ is the initial state distribution. We denote the trajectory of an episode by $\tau = (s_0, a_0, r_0, \dots, s_T, a_T, r_T, s_{T+1})$ where $r_t = R(s_t, a_t)$ and $T$ is the length of the trajectory which can be infinite. 
%
If a trajectory is generated by a stochastic policy $\pi: \gS \times \gA \rightarrow \sR_{\geq 0}$, $R^\pi = \sum_{t=0}^T \gamma^t r_t$ is a random variable that describes the \textit{discounted return} the policy achieves. The objective is to find a policy $\pi^\star$ that maximizes the expected discounted return in the MDP.
%$\pi^\star = \argmax_\pi \E_{\tau \sim p^\pi(\cdot)}\left[ Z^\pi\right], \,\, \text{where} \,\, p^\pi(\tau) = \mu(\vs_0) \prod_{t=0}^T P(\vs_{t+1} \mid \vs_{t}, \va_{t}) \pi(\vs_{t}\mid \va_{t}).$ 

\textbf{Problem Setting.}
In this paper, we are interested in learning new tasks from a small number of demonstrations after pretraining on a large dataset of offline trajectories from different tasks. Formally, we consider a set of tasks $\mathcal{T}$ split into disjoint sets $\mathcal{T}_{train}$ and $\mathcal{T}_{test}$.
%We consider $N$ training and $M$ test tasks, $\mathcal{T}_{train} \in \mathcal{T}$ and $\mathcal{T}_{test} \in \mathcal{T}$, respectively, such that $\mathcal{T}_{train} \cap \mathcal{T}_{train} = \emptyset $. 
The train and test tasks have the same action space, but different state spaces, initial state distributions, transition and reward functions. Each task $\mathcal{T}_i \in \mathcal{T}$ is procedurally generated, meaning that it includes $L$ MDPs $\{\mathcal{M}_i^l\}_{l=1}^L$ sharing the same state spaces, transition and reward functions, but different initial state distributions. We refer to these as ``levels" in this paper as depicted in \Cref{fig:train_test}: for MiniHack and Procgen, each level corresponds to a different random seed used to generate the corresponding instance of the task. For each MDP $\mathcal{M}_i$ in our set, we collect $K$ expert demonstrations using reinforcement learning (RL) agents which were pretrained on $\mathcal{M}_i$, to create a dataset $\mathcal{D}_i$ containing trajectories of the form $\tau_i = (s_0^i, a_0^i, \ldots, s_T^i, a_T^i)$, where $T$ is the maximum number of steps in the episode. For the test MDPs $K$ is small, meaning that we only collect a few expert demonstrations for each level of a test task and use them to condition our transformer model. 
 
